% R. Brown's translations of two letters to L. Breen, with his note and comments (cote 134-2): archive cote 134-2, site page 120 (PDF page 121).
% comments. Draft compared: none.
% AI-transcribed by Claude Opus 5.5 (claude-opus-5-5, adaptive thinking, effort high), Message Batches API, batch brown_full (msgbatch_01Md3i5yK1PxkpJn5GHfEExE), 2026-11-13; not yet checked by eye.
\scanpage{134-2/120}
\mspage{-4-}

The latter paper uses the equivalence between crossed complexes
and simplicial T-complexes (which generalises the Dold-Kan
theorem on the equivalence of chain complexes and simplicial
abelian groups) to give a notion of nerve, and so classifying
space BC of a crossed complex C. The classical Eilenberg-%
MacLane bijection
\[ [X, K(A, n)] \cong H^{n} (X, A) \]
for a CW-complex X, is then generalised to a bijection
\[ [X, BC] \cong H^{0} (X, C) \]
However, as pointed out in (5), these results are still ``too
abelian'', and much more work is needed to realise Grothendieck's
programme.

10. The theory of tubular neighbourhoods has been developed in papers
by D.A. Cox.

'Algebraic tubular neighbourhoods I, II', \unsure{Math,|Math.} Scand. 42 (1978),
211 - 228, 229 - 242.

Another of his papers seems relevant:

'Homotopy theory of simplicial schemes', Compositio Math. 39(1979)
263 - 296.

R. BROWN

School of Mathematics and Computer Science,
University College of North Wales,
Bangor,
Gwynedd. \hand{U.K.}

November, 1982

\margin{11 A van Kampen Theorem for 2-cat-groups has now
been proved by Brown and Loday.}

\mspage{\hand{121}}
